\chapter{2007年山东卷（文）}

\section{单选题}

\begin{problem}
复数 \(\frac{4 + 3\i}{1 + 2\i}\) 的实部是（\quad）

\choices
    {\(-2\)}
    {\(2\)}
    {\(3\)}
    {\(4\)}
\end{problem}

\begin{answer}
B.
\end{answer}

\begin{problem}
已知集合 \(M = \{-1, 1\}\)，\(N = \left\{x \mid \frac{1}{2} < 2^{x+1} < 4, x \in \Z\right\}\)，则 \(M \cap N =\)（\quad）

\choices
    {\( \{-1,1\} \)}
    {\( \{-1\} \)}
    {\( \{0\} \)}
    {\( \{-1,0\} \)}
\end{problem}

\begin{answer}
B.
\end{answer}

\begin{problem}
下列几何体各自的三视图中，有且仅有两个视图相同的是（\quad）

\bitmapfigure[max width=0.84\linewidth]{img/2007/shandong_liberal/q3_fig1.png}

\choices
    {\circled{1}\circled{2}}
    {\circled{1}\circled{3}}
    {\circled{1}\circled{4}}
    {\circled{2}\circled{4}}
\end{problem}

\begin{answer}
D.
\end{answer}

\begin{problem}
要得到函数 \( y = \sin x \) 的图象，只需将函数 \( y = \cos \left(x - \frac{\pi}{3}\right) \) 的图象（\quad）

\choices
    {向右平移 \(\frac{\pi}{6}\) 个单位}
    {向右平移 \(\frac{\pi}{3}\) 个单位}
    {向左平移 \(\frac{\pi}{3}\) 个单位}
    {向左平移 \(\frac{\pi}{6}\) 个单位}
\end{problem}

\begin{answer}
A.
\end{answer}

\begin{problem}
已知向量 \(\bs{a} = (1, n)\)，\(\bs{b} = (-1, n)\)，若 \(2\bs{a} - \bs{b}\) 与 \(\bs{b}\) 垂直，则 \(|\bs{a}| =\)（\quad）

\choices
    {\(1\)}
    {\(\sqrt{2}\)}
    {\(2\)}
    {\(4\)}
\end{problem}

\begin{answer}
C.
\end{answer}

\begin{problem}
给出下列三个等式：\(f(xy) = f(x) + f(y)\)，\(f(x + y) = f(x)f(y)\)，\(f(x + y) = \frac{f(x) + f(y)}{1 - f(x)f(y)}\). 下列函数中不满足其中任何一个等式的是（\quad）

\choices
    {\(f(x) = 3^{x}\)}
    {\(f(x) = \sin x\)}
    {\(f(x) = \log_2x\)}
    {\(f(x) = \tan x\)}
\end{problem}

\begin{answer}
B.
\end{answer}

\begin{problem}
命题“对任意的 \(x \in \R\)，\(x^3 - x^2 + 1 \leqslant 0\)”的否定是（\quad）

\choices
    {不存在 \(x \in \R, x^3 - x^2 + 1 \leqslant 0\)}
    {存在 \(x \in \R\)，\(x^3 - x^2 + 1 \leqslant 0\)}
    {存在 \(x \in \R\)，\(x^3 - x^2 + 1 > 0\)}
    {对任意的 \(x \in \R\)，\(x^3 - x^2 + 1 > 0\)}
\end{problem}

\begin{answer}
C.
\end{answer}

\begin{problem}
某班 \(50\) 名学生在一次百米测试中，成绩全部介于 \(13\) 秒与 \(19\) 秒之间，将测试结果按如下方式分成六组：第一组，成绩大于等于 \(13\) 秒且小于 \(14\) 秒；第二组，成绩大于等于 \(14\) 秒且小于 \(15\) 秒；…；第六组，成绩大于等于 \(18\) 秒且小于 \(19\) 秒. 如图是按上述分组方法得到的频率分布直方图. 设成绩小于 \(17\) 秒的学生人数占全班总人数的百分比为 \(x\)，成绩大于等于 \(15\) 秒且小于 \(17\) 秒的学生人数为 \(y\)，则从频率分布直方图中可分析出 \(x\) 和 \(y\) 分别为（\quad）

\texfigure[max width=0.46\linewidth]{img_repaint/2007/shandong_liberal/q8_fig1.tex}

\choices
    {\(0.9, 35\)}
    {\(0.9, 45\)}
    {\(0.1, 35\)}
    {\(0.1, 45\)}
\end{problem}

\begin{answer}
A.
\end{answer}

\begin{problem}
设 \(O\) 是坐标原点，\(F\) 是抛物线 \( y^{2} = 2px \)（\(p > 0\)）的焦点，\(A\) 是抛物线上的一点，\( \overrightarrow{FA} \) 与 \(x\) 轴正向的夹角为 \( 60^{\circ} \)，则 \( \left|\overrightarrow{OA}\right| \) 为（\quad）

\choices
    {\(\frac{21p}{4}\)}
    {\(\frac{\sqrt{21}p}{2}\)}
    {\(\frac{\sqrt{13}}{6} p\)}
    {\(\frac{13}{36} p\)}
\end{problem}

\begin{answer}
B.
\end{answer}

\begin{problem}
阅读如图的程序框图，若输入的 \(n\) 是 \(100\)，则输出的变量 \(S\) 和 \(T\) 的值依次是（\quad）

\texfigure[max width=0.55\linewidth]{img_repaint/2007/shandong_liberal/q10_fig1.tex}

\choices
    {\(2550, 2500\)}
    {\(2550, 2550\)}
    {\(2500, 2500\)}
    {\(2500, 2550\)}
\end{problem}

\begin{answer}
A.
\end{answer}

\begin{problem}
设函数 \( y = x^3 \) 与 \( y = \left(\frac{1}{2}\right)^{x - 2} \) 的图象的交点为 \( (x_0, y_0) \)，则 \( x_0 \) 所在的区间是（\quad）

\choices
    {\((0,1)\)}
    {\((1,2)\)}
    {\((2,3)\)}
    {\((3,4)\)}
\end{problem}

\begin{answer}
B.
\end{answer}

\begin{problem}
设集合 \(A = \{1, 2\}\)，\(B = \{1, 2, 3\}\)，分别从集合 \(A\) 和 \(B\) 中随机取一个数 \(a\) 和 \(b\)，确定平面上的一个点 \(P(a, b)\). 记“点 \(P(a, b)\) 落在直线 \(x + y = n\) 上”为事件 \(C_n\) （\(2 \leqslant n \leqslant 5\)，\(n \in \N\)）. 若事件 \(C_n\) 的概率最大，则 \(n\) 的所有可能值为（\quad）

\choices
    {\(3\)}
    {\(4\)}
    {\(2\) 和 \(5\)}
    {\(3\) 和 \(4\)}
\end{problem}

\begin{answer}
D.
\end{answer}

\section{填空题}

\begin{problem}
设函数 \( f_{1}(x)=x^{1/2} \)，\( f_{2}(x)=x^{-1} \)，\( f_{3}(x)=x^{2} \)，则 \( f_{1}(f_{2}(f_{3}(2007))) = \)\fillinblank{}.
\end{problem}

\begin{answer}
\(\frac{1}{2007}\)
\end{answer}

\begin{problem}
函数 \( y = a^{1 - x} \)（\(a > 0\)，\(a \neq 1\)）的图象恒过定点 \(A\)，若点 \(A\) 在直线 \( mx + ny - 1 = 0 \)（\(mn > 0\)）上，则 \( \frac{1}{m} + \frac{1}{n} \) 的最小值为\fillinblank{}.
\end{problem}

\begin{answer}
\(4\)
\end{answer}

\begin{problem}
当 \(x \in (1,2)\) 时，不等式 \(x^2 + mx + 4 < 0\) 恒成立，则 \(m\) 的取值范围是\fillinblank{}.
\end{problem}

\begin{answer}
\(m\leqslant-5\)
\end{answer}

\begin{problem}
与直线 \(x + y - 2 = 0\) 和曲线 \(x^{2} + y^{2} - 12x - 12y + 54 = 0\) 都相切的半径最小的圆的标准方程是\fillinblank{}.
\end{problem}

\begin{answer}
\((x-2)^2+(y-2)^2=2\)
\end{answer}

\section{解答题}

\begin{problem}
在 \(\triangle ABC\) 中，角 \(A\)，\(B\)，\(C\) 的对边分别为 \(a\)，\(b\)，\(c\)，\(\tan C = 3\sqrt{7}\).

\begin{enumerate}
\item 求 \( \cos C \).
\item 若 \( \overrightarrow{CB} \cdot \overrightarrow{CA} = \frac{5}{2} \)，且 \( a + b = 9 \)，求 \(c\).
\end{enumerate}
\end{problem}

\begin{solution}
\begin{enumerate}
\item 由 \(0<C<\pi\) 且 \(\tan C=3\sqrt{7}>0\)，知 \(C\) 为锐角，故 \(\cos C>0\). 因为 \(1+\tan^{2}C=1+\left(3\sqrt{7}\right)^{2}=64\)，所以 \(\cos C=\dfrac{1}{\sqrt{64}}=\dfrac{1}{8}\).
\item 由向量数量积公式，得 \(\overrightarrow{CB}\cdot\overrightarrow{CA}=ab\cos C\)，从而 \(\dfrac{ab}{8}=\dfrac{5}{2}\)，即 \(ab=20\). 又因为 \(a+b=9\)，所以由余弦定理，
\[
c^{2}=a^{2}+b^{2}-2ab\cos C=(a+b)^{2}-2ab-2ab\cos C=81-40-5=36.
\]
由于 \(c>0\)，故 \(c=6\).
\end{enumerate}
\end{solution}

\begin{problem}
设 \(\{a_{n}\}\) 是公比大于 \(1\) 的等比数列，\(S_{n}\) 为数列 \(\{a_{n}\}\) 的前 \(n\) 项和. 已知 \(S_{3} = 7\)，且 \(a_{1} + 3\)，\(3a_{2}\)，\(a_{3} + 4\) 构成等差数列.

\begin{enumerate}
\item 求数列 \( \{a_{n}\} \) 的通项；
\item 令 \(b_{n} = \ln a_{3n + 1}\)，\(n = 1\)，\(2\)，\(\cdots\)，求数列 \( \{b_{n}\} \) 的前 \( n \) 项和 \( T_{n} \).
\end{enumerate}
\end{problem}

\begin{solution}
\begin{enumerate}
\item 设等比数列的公比为 \(q\)，则 \(q>1\)，且 \(a_{1}\left(1+q+q^{2}\right)=7\). 由三个数构成等差数列，得
\[
2\cdot 3a_{2}=(a_{1}+3)+(a_{3}+4),
\]
即 \(6a_{1}q=a_{1}\left(1+q^{2}\right)+7\). 由 \(a_{1}\left(1+q^{2}\right)=7-a_{1}q\)，得 \(6a_{1}q=14-a_{1}q\)，所以 \(a_{1}q=2\)，即 \(a_{2}=2\).

又 \(a_{1}+a_{3}=S_{3}-a_{2}=5\)，且等比数列满足 \(a_{1}a_{3}=a_{2}^{2}=4\)，因此 \(a_{1}\)，\(a_{3}\) 是方程 \(t^{2}-5t+4=0\) 的两个根，即 \(\{a_{1},a_{3}\}=\{1,4\}\). 因为 \(q>1\)，只能取 \(a_{1}=1\)，\(a_{3}=4\)，从而 \(q=2\). 所以数列的通项为 \(a_{n}=2^{n-1}\).
\item 由 \(a_{n}=2^{n-1}\)，得 \(a_{3n+1}=2^{3n}\)，所以 \(b_{n}=\ln a_{3n+1}=3n\ln 2\). 因而
\[
T_{n}=\sum_{k=1}^{n}b_{k}=3\ln 2\sum_{k=1}^{n}k=\frac{3n(n+1)}{2}\ln 2.
\]
\end{enumerate}
\end{solution}

\begin{problem}
本公司计划 \(2008\) 年在甲，乙两个电视台做总时间不超过 \(300\) 分钟的广告，广告总费用不超过 \(9\) 万元. 甲，乙电视台的广告收费标准分别为 \(500\) 元/分钟和 \(200\) 元/分钟. 规定甲，乙两个电视台为该公司所做的每分钟广告，能给公司带来的收益分别为 \(0.3\) 万元和 \(0.2\) 万元. 问该公司如何分配在甲，乙两个电视台的广告时间，才能使公司的收益最大，最大收益是多少万元？
\end{problem}

\begin{solution}
设在甲、乙两个电视台分别投放广告 \(x\) 分钟、\(y\) 分钟，则 \(x\geqslant0\)，\(y\geqslant0\)，并且费用单位统一为万元后，约束条件为
\[
\begin{cases}
x+y\leqslant300,\\
0.05x+0.02y\leqslant9.
\end{cases}
\]
公司的收益为 \(R=0.3x+0.2y\). 可行域的顶点为 \((0,0)\)，\((0,300)\)，\((180,0)\) 以及两条约束边界的交点. 联立 \(x+y=300\) 与 \(0.05x+0.02y=9\)，得交点为 \((100,200)\).

分别计算各顶点的收益：\(R(0,0)=0\)，\(R(0,300)=60\)，\(R(180,0)=54\)，\(R(100,200)=70\). 因此，当 \(x=100\)，\(y=200\) 时收益最大，最大收益为 \(70\) 万元.
\end{solution}

\begin{problem}
如图，在直四棱柱 \(ABCD - A_{1}B_{1}C_{1}D_{1}\) 中，已知 \(DC = DD_{1} = 2AD = 2AB\)，\(AD \perp DC\)，\(AB \parallel DC\).

\begin{enumerate}
\item 求证：\(D_{1}C \perp AC_{1}\).
\item 设 \(E\) 是 \(DC\) 上一点，试确定 \(E\) 的位置，使 \(D_{1}E \parallel\) 平面 \(A_{1}BD\)，并说明理由.
\end{enumerate}

\bitmapfigure[max width=0.36\linewidth]{img/2007/shandong_liberal/q20_fig1.png}
\end{problem}

\begin{solution}
设 \(AD=AB=t>0\). 建立空间直角坐标系，使底面 \(ABCD\) 在 \(xOy\) 平面内，并取 \(D=(0,0,0)\)，\(C=(2t,0,0)\)，\(A=(0,t,0)\)，\(B=(t,t,0)\). 由 \(DD_{1}=2t\)，得 \(D_{1}=(0,0,2t)\)，\(A_{1}=(0,t,2t)\)，\(C_{1}=(2t,0,2t)\).

\begin{enumerate}
\item \(\overrightarrow{D_{1}C}=(2t,0,-2t)\)，\(\overrightarrow{AC_{1}}=(2t,-t,2t)\)，所以
\[
\overrightarrow{D_{1}C}\cdot\overrightarrow{AC_{1}}=2t\cdot2t+0\cdot(-t)+(-2t)\cdot2t=0.
\]
故 \(D_{1}C\perp AC_{1}\).
\item 设 \(E=(s,0,0)\)，其中 \(0\leqslant s\leqslant2t\). 平面 \(A_{1}BD\) 内的两个不共线向量为 \(\overrightarrow{DA_{1}}=(0,t,2t)\) 和 \(\overrightarrow{DB}=(t,t,0)\)，故其一个法向量为
\[
\bs n=\overrightarrow{DA_{1}}\times\overrightarrow{DB}=t^{2}(-2,2,-1).
\]
又 \(\overrightarrow{D_{1}E}=(s,0,-2t)\). 要使 \(D_{1}E\parallel\) 平面 \(A_{1}BD\)，其方向向量应与该平面的法向量垂直，于是
\[
\overrightarrow{D_{1}E}\cdot\bs n=2t^{2}(t-s)=0.
\]
由于 \(t>0\)，得 \(s=t\). 此时 \(E=(t,0,0)\)，即 \(E\) 为 \(DC\) 的中点. 平面 \(A_{1}BD\) 的方程为 \(-2x+2y-z=0\)，而 \(D_{1}\) 不在此平面内；因此当 \(s=t\) 时，直线 \(D_{1}E\) 与该平面无公共点且方向平行于该平面，故 \(D_{1}E\parallel\) 平面 \(A_{1}BD\).
\end{enumerate}
\end{solution}

\begin{problem}
设函数 \( f(x) = ax^2 + b\ln x \)，其中 \( ab \neq 0 \). 证明：当 \( ab > 0 \) 时，函数 \( f(x) \) 没有极值点；当 \( ab < 0 \) 时，函数 \( f(x) \) 有且只有一个极值点，并求出极值.
\end{problem}

\begin{solution}
函数的定义域为 \(x>0\)，且
\[
f'(x)=2ax+\frac{b}{x}=\frac{2ax^{2}+b}{x}.
\]
当 \(ab>0\) 时，\(a\) 与 \(b\) 同号，所以对任意 \(x>0\)，\(2ax^{2}+b\) 不为零，进而 \(f'(x)\) 不为零，函数没有极值点.

当 \(ab<0\) 时，方程 \(f'(x)=0\) 等价于 \(2ax^{2}+b=0\)，在定义域内有且只有一个解
\[
x_{0}=\sqrt{-\frac{b}{2a}}.
\]
又
\[
f''(x)=2a-\frac{b}{x^{2}},\qquad f''(x_{0})=4a\neq0,
\]
所以 \(x_{0}\) 是唯一的极值点；当 \(a>0\) 时为极小值点，当 \(a<0\) 时为极大值点. 极值为
\[
f(x_{0})=a\left(-\frac{b}{2a}\right)+b\ln\sqrt{-\frac{b}{2a}}=\frac{b}{2}\left(\ln\left(-\frac{b}{2a}\right)-1\right).
\]
\end{solution}

\begin{problem}
已知椭圆 \(C\) 的中心在坐标原点，焦点在 \(x\) 轴上，椭圆 \(C\) 上的点到焦点距离的最大值为 \(3\)，最小值为 \(1\).

\begin{enumerate}
\item 求椭圆 \(C\) 的标准方程.
\item 若直线 \(l: y = kx + m\) 与椭圆 \(C\) 相交于 \(A\)，\(B\) 两点（\(A\)，\(B\) 不是左右顶点），且以 \(AB\) 为直径的圆过椭圆 \(C\) 的右顶点，求证：直线 \(l\) 过定点，并求出该定点的坐标.
\end{enumerate}
\end{problem}

\begin{solution}
\begin{enumerate}
\item 设椭圆的长半轴、短半轴和半焦距分别为 \(a\)，\(b\)，\(c\)，其中 \(a>b>0\)，则 \(c^{2}=a^{2}-b^{2}\). 对于右焦点，椭圆上点到它的最大距离和最小距离分别为 \(a+c\) 和 \(a-c\)，所以
\[
a+c=3,\qquad a-c=1.
\]
解得 \(a=2\)，\(c=1\)，进而 \(b^{2}=a^{2}-c^{2}=3\). 因此椭圆的标准方程为
\[
\frac{x^{2}}{4}+\frac{y^{2}}{3}=1.
\]
\item 设直线 \(l\) 与椭圆的两个交点横坐标为 \(x_{1}\)，\(x_{2}\)，则对应纵坐标为 \(y_{i}=kx_{i}+m\)（\(i=1,2\)）. 把直线方程代入椭圆方程，得
\[
3x^{2}+4(kx+m)^{2}=12.
\]
由韦达定理，
\[
x_{1}+x_{2}=-\frac{8km}{3+4k^{2}},\qquad x_{1}x_{2}=\frac{4m^{2}-12}{3+4k^{2}}.
\]
椭圆的右顶点为 \(R=(2,0)\). 以 \(AB\) 为直径的圆过 \(R\)，等价于 \(\overrightarrow{RA}\cdot\overrightarrow{RB}=0\)，即
\[
(2-x_{1})(2-x_{2})+y_{1}y_{2}=0.
\]
将 \(y_{i}=kx_{i}+m\) 代入并整理，得
\[
4+(km-2)(x_{1}+x_{2})+(1+k^{2})x_{1}x_{2}+m^{2}=0.
\]
再代入韦达定理并化简，得到
\[
4k^{2}+16km+7m^{2}=0,
\]
即
\[
(2k+m)(2k+7m)=0.
\]
若 \(2k+m=0\)，则 \(l\) 的方程为 \(y=k(x-2)\)，直线经过右顶点 \(R\)，这与 \(A\)，\(B\) 不是左右顶点矛盾. 因而必有 \(2k+7m=0\)，从而
\[
l:y=k\left(x-\frac{2}{7}\right).
\]
所以所有符合条件的直线均经过定点 \(\left(\frac{2}{7},0\right)\).
\end{enumerate}
\end{solution}
